\documentclass[12pt]{article}
\usepackage[T1]{fontenc}
\usepackage[utf8]{inputenc}
\usepackage{lmodern}
\usepackage{textcomp}
\usepackage{enumitem}
\usepackage{geometry}
\geometry{margin=1in}
\begin{document}
\begin{center}
Answer Key - Philosophy - Graduate Level Assessment 2 \
\small (Answers are ordered exactly as in the question paper.)
\end{center}

\section*{Multiple Choice}
\begin{enumerate}[label=\arabic*.]
\item \textbf{Answer:} A
\\ \textit{Options:} A) The Bees do not win their first game.; B) Both the Aardvarks and the Chipmunks do not win their first games.; C) The Chipmunks do not win their first game.; D) The Chipmunks win their first game.; E) Both the Aardvarks and the Chipmunks win their first games.
\\ \textit{Note:} The correct answer is "The Bees do not win their first game" because, in the conditional proposition "A only if B," the antecedent is the condition that must be true for the consequent to hold, which, in this case, indicates that for the Bees to win their game, it is necessary that either the Aardvarks or the Chipmunks do not win theirs.
\item \textbf{Answer:} A
\\ \textit{Options:} A) Discipline and order; B) Strength and perseverance; C) Wisdom and knowledge; D) Goodwill and love; E) Courage and bravery
\\ \textit{Note:} Guru Nanak and his successors emphasized the use of specific Sikh Ragas to instill a sense of discipline and order within the practice of worship and community life, thereby promoting a structured approach to spirituality and devotion in Sikhism.
\item \textbf{Answer:} A
\\ \textit{Options:} A) John's duty to return to Mary that car that he borrowed from her; B) John's duty to allow Mary to pursue goals that she values; C) John's duty to not harm Mary; D) John's duty to not commit suicide
\\ \textit{Note:} This is the correct answer because John's duty to return the borrowed car to Mary exemplifies the reciprocal relationship between rights and duties, where Mary's right to her property (the car) necessitates John's obligation to return it.
\item \textbf{Answer:} A
\\ \textit{Options:} A) In faith in God; B) In the strict adherence to the Five Ks; C) In ascetic renunciation; D) In the practice of yoga and other physical disciplines; E) In acts of charity and service
\\ \textit{Note:} The key to liberation in Sikhism is found in faith in God because it emphasizes the belief that surrendering to the divine will and cultivating a deep spiritual connection with God leads to liberation from the cycle of birth and death.
\item \textbf{Answer:} D
\\ \textit{Options:} A) the postulate of immortality.; B) the autonomous will.; C) the categorical imperative.; D) legislating for oneself.; E) the unity of consciousness.
\\ \textit{Note:} Anscombe criticizes Kant's idea of legislating for oneself because she argues that it undermines the objective nature of moral law, suggesting that self-legislation leads to a relativistic view of ethics that is inconsistent with the notion of universal moral principles.
\end{enumerate}

\section*{True / False}
\begin{enumerate}[label=\arabic*.]
\setcounter{enumi}{5}
\item \textbf{Answer:} True
\\ \textit{Options:} A) True; B) False
\\ \textit{Note:} The statement is true because the execution of innocent persons is often discussed as a significant ethical concern against the death penalty, indicating that it is indeed recognized as a potential negative consequence, rather than being dismissed as irrelevant.
\item \textbf{Answer:} False
\\ \textit{Options:} A) True; B) False
\\ \textit{Note:} This statement is false because John Locke is primarily known for his contributions to social contract theory and empiricism, rather than a version of virtue ethics that emphasizes moral character and virtue cultivation.
\item \textbf{Answer:} True
\\ \textit{Options:} A) True; B) False
\\ \textit{Note:} Rawls' theory comprises the principles of justice and the original position, each of which can be understood and accepted individually, demonstrating their conceptual independence within his framework.
\item \textbf{Answer:} False
\\ \textit{Options:} A) True; B) False
\\ \textit{Note:} The ambrosial hours for Sikhs, known as "Amrit Vela," are traditionally observed from approximately 3 to 6 a.m., not from 1 to 4 a.m., emphasizing the importance of early morning meditation and worship.
\item \textbf{Answer:} True
\\ \textit{Options:} A) True; B) False
\\ \textit{Note:} The moral argument against imprisoning debtors is true because the principle of universalizability asserts that moral principles should apply consistently to all individuals in similar circumstances, highlighting the inconsistency and unfairness of punishing individuals for failing to meet financial obligations that may be beyond their control.
\end{enumerate}

\section*{Long Form Response}
\begin{enumerate}[label=\arabic*.]
\setcounter{enumi}{10}
\item \textbf{Answer:} The term "fearful" within Hindu traditions carries profound philosophical implications that significantly influence gender roles and societal perceptions of women. Traditionally, women are often portrayed as embodiments of fear, particularly in their relationship to male authority and societal expectations, which can lead to their subjugation and limit their autonomy. This characterization reinforces a perception of women as inherently vulnerable, which in turn perpetuates a cycle of dependency and reinforces patriarchal structures. Furthermore, the association of women with fear can undermine their agency, positioning them as passive recipients of societal norms rather than active participants in their own lives. Ultimately, examining the implications of this term reveals how philosophical constructs shape and constrain women's identities and roles within both religious and social contexts.
\\ \textit{Note:} The answer correctly identifies how the portrayal of women as "fearful" in Hindu traditions shapes their gender roles and societal perceptions by emphasizing their subjugation, vulnerability, and dependency on male authority, thus reflecting broader philosophical implications regarding autonomy and power dynamics.
\item \textbf{Answer:} G.E. Moore's philosophical framework is anchored in several key concepts, including the naturalistic fallacy, the open question argument, and his defense of common sense realism. The naturalistic fallacy challenges the conflation of moral properties with natural properties, asserting that "good" cannot be defined through empirical means, as exemplified in his famous phrase "good is good." The open question argument further emphasizes this by illustrating that any attempt to define moral terms can be questioned, indicating that moral truths are distinct from factual truths. Additionally, Moore's commitment to common sense realism underscores his belief in the existence of an external reality independent of human perception, which interrelates with his ethical theories by grounding moral claims in a reality that we can recognize through intuition. Collectively, these concepts form a cohesive philosophical stance that underscores the complexity and depth of Moore's contributions to ethics and epistemology.
\\ \textit{Note:} correct as it accurately identifies and elaborates on G.E. Moore's central concepts-the naturalistic fallacy, the open question argument, and common sense realism-demonstrating their significance in his philosophical framework and their interrelatedness in the context of moral philosophy.
\end{enumerate}

\vfill
\noindent\textit{End of Answer Key}
\end{document}